\clearpage
\section{Einleitung}
\label{sec:einleitung}

Die Gemeinde Schwanden bei Brienz im Kanton Bern liegt am Fuss eines aktiven Bergsturzgebiets
an der Südflanke des Brienzer Rothorns. Nach dem letzten grossen Bergsturz von 1901
wurden verschiedene Schutzmassnahmen umgesetzt. Neue Bewegungen im Bereich der Ägerdi
Spalte zeigten in den 1980er-Jahren jedoch, dass der Hang weiterhin in Bewegung ist.
Seit 1989 überwacht das Institut Geomatik der FHNW im Auftrag der Gemeinde den Hang
mit millimetergenauen geodätischen Messungen. Die so über Jahrzehnte aufgebaute 
Messreihe ist eine wichtige Grundlage, um das Bewegungsverhalten zu beurteilen und 
die Bevölkerung sowie die darunterliegenden Siedlungen zu schützen.


Im Herbstsemester 2026 wurde die Überwachungsmessung im Feldkurs des Vertiefungsprofils
GeoSensorik und Monitoring erneut durchgeführt. Im anschliessenden Modul Geosensorik und 
Monitoring II (5220) werden die erhobenen Daten aufbereitet und ausgewertet. Das Ziel 
ist, die Messreihe mit einer neuen, zuverlässigen Epoche zu verlängern. Zudem sollen 
die Auswertungen als Entscheidungsgrundlage dienen, die Messanordnung weiter zu vereinfachen
und lange terrestrische Beobachtungen weitgehend durch langstatische GNSS-Messungen 
zu ersetzen

\subsection{Ausgangslage und Auftrag}
Mit den im Feldkurs 2026 erhobenen Messdaten soll die langjährige Messreihe am
Schwanderbärgli verlängert werden. Dazu sind eine sorgfältige Aufbereitung aller Daten
und eine zielführende Netzausgleichung nötig. Die Messdaten stammen aus verschiedenen
Sensoren (Totalstation, GNSS, Nivellement, trigonometrische Höhenbestimmung) und müssen
zusammengeführt sowie auf Vollständigkeit und Konsistenz geprüft werden. Das Ergebnis
ist ein zuverlässiger Koordinatensatz mit bekannter Qualität, der die Grundlage für
den Epochenvergleich und die Deformationsanalyse bildet.

\subsection{Ziele und Abgrenzung}
Die Auswertung der Felddaten ist auf mehrere Gruppen aufgeteilt, die jeweils einen
Teil bearbeiten. Unsere Gruppe hat zwei Arbeitspakete übernommen:
\begin{enumerate}
    \item \textbf{Datenaufbereitung der tachymetrischen Messungen (TPS):} Import,
    Strukturierung und Vereinheitlichung der Messdaten, Kontrolle auf Vollständigkeit
    und Plausibilität, Anbringen der erforderlichen Korrekturen sowie Festlegung,
    welche Beobachtungen für die Ausgleichung verwendet werden.
    \item \textbf{Auswertung des Netzes Schwanderbärgli in der Variante V1.1
    «optimiert»:} Netzausgleichung nach der Methode der kleinsten Quadrate mit einer
    von uns vorgeschlagenen optimierten Konfiguration, inklusive Beurteilung von
    Genauigkeit und Zuverlässigkeit.
\end{enumerate}


